\section{Love, Freedom, and Joy Without End}

The vision of God is inseparable from loving him and rejoicing in him. We can know a created good and remain divided about it, fearing the cost of receiving it or preferring some competing satisfaction. In the blessed, charity is whole. God is known as the supreme good and loved without reserve. Their joy is the joy of possessing him, and their security belongs to the same gift. The fear that makes earthly happiness tremble is gone. They do not watch the good recede, calculate how long it may last, or brace themselves against the next separation.

Augustine ends \emph{City of God} by considering this secure happiness (XXII.30). He describes the blessed resting, seeing, loving, and praising. The blessed are free from the weariness and need that burden our action. Their praise expresses the goodness they behold and the gratitude of those who have received everything. God himself is the reward. Augustine gathers the rightful objects of human longing---life, health, peace, honor, abundance---into the possession of the giver. Their fulfillment rests in him, and their many forms of goodness no longer pull the heart apart.

Eternity cannot be understood simply by extending the duration of an earthly pleasure. Our pleasures can become tedious because they are limited, because our powers tire, or because desire seeks a good they cannot supply. These conditions do not define the blessed enjoyment of God. Augustine speaks of a love without satiety in the sense of weariness, and praise without exhaustion. The blessed continue to delight in the God whom they possess. We cannot describe the entire manner of a creature's eternal experience, but we can entrust ourselves to the promise that the good possessed will never become a burden to those who possess it.

The freedom of the blessed belongs within this fullness. Augustine says that the final freedom is greater than the first: the blessed cannot sin, and their inability is God's gift. Their will has been healed so completely that delight in evil has no place in it. To choose sin is to damage the freedom for which we were made. The power to turn against our good is therefore no perfection whose absence would impoverish heaven. God himself cannot sin; the creature shares, by grace, an unfailing adherence to goodness without becoming divine by nature.

This teaching reaches the weariness of repeated moral struggle. We know the divided experience of wanting to be generous and resenting the request, desiring purity and returning to what harms us, intending to forgive and rehearsing the offense. Grace already works within those conflicts, but the conflict is not our final condition. Heaven promises a will no longer at war with itself. The good will be loved with the whole person. Perseverance gains strength from knowing that the present battle is directed toward peace, and that peace will involve a more complete exercise of love than we have yet known.

The desire for such freedom changes how we practice virtue. A reluctant act of patience may be a true beginning, yet we hope to become patient enough to delight in another person's good. We resist a temptation, but we also seek the healing of the desire that makes it powerful. We ask forgiveness for an outburst, then ask to love the person whose need exposed our impatience. These acts seek an inward conversion that restraint alone cannot supply. We ask God to heal the desires behind our conduct, until we can love the good willingly and without division.

Paul's account of charity places this fulfillment before us with great sobriety (1~Cor 13). Prophetic gifts, eloquence, knowledge, and sacrifice can dazzle observers while the heart remains without love. Charity is patient, kind, free from envy, and joined to truth. Its permanence gives the passing duties of our lives their eternal seriousness. We shall not carry every earthly capacity into glory in its present use, but love formed by grace belongs to the life that endures. Paul gives us reason to examine an admired achievement by the same measure as a hidden act of service: the charity with which it was done.

The annual remembrance of heaven can therefore become a renewal of courage. God is drawing us toward a happiness in which we can love without the shame of divided intention, rejoice without the fear of theft or loss, and praise without turning praise back upon ourselves. A habit that repeatedly defeats us can still be brought to him for healing. We continue asking because he promises a freedom more complete than the self-restraint we have so far attained.

